% !PRINTED 12
\weberpage{12}

Setzen wir, wie in § 2

\weq{5}{\begin{gathered}
a = (x_2x_3)(x_1x_4), \\
b = (x_3x_1)(x_2x_4), \\
c = (x_1x_2)(x_3x_4), \\
a + b + c = 0,
\end{gathered}}

so könnnen wir, da es nur auf das Verhältnis von $\alpha$ zu $\beta$ ankommt, die letzte der Gleichungen (4) dadurch befriedigen, daß wir setzen: % [sic] printed with three n-strokes

\weq{6}{\sqrt{\alpha} = \sqrt{\pm a} + \sqrt{\mp b}, \quad \sqrt{\beta} = \sqrt{\pm a} - \sqrt{\mp b},}

und aus (3) ergibt sich für $z$ der Ausdruck

\weq{7}{z = \sqrt{\alpha}\,\dfrac{(xx_2)\sqrt{(x_3x_1)(x_4x_1)} + (xx_1)\sqrt{(x_3x_2)(x_4x_2)}}{(xx_2)\sqrt{(x_3x_1)(x_4x_1)} - (xx_1)\sqrt{(x_3x_2)(x_4x_2)}}.}

Stellt man neben (3) noch die daraus folgende Gleichung:

\[
h'\,\dfrac{x-x_3}{x-x_4} = \dfrac{z-\sqrt{\beta}}{z+\sqrt{\beta}}
\]

auf, so ergibt sich durch logarithmische Differentiation:

\[
\dfrac{(x_1x_2)\,dx}{(xx_1)(xx_2)} = \dfrac{2\sqrt{\alpha}\,dz}{z^2-\alpha}, \quad \dfrac{(x_3x_4)\,dx}{(xx_3)(xx_4)} = \dfrac{2\sqrt{\beta}\,dz}{z^2-\beta},
\]

und daraus durch Multiplikation mit Rücksicht auf (5) und (6):

\weq{8}{\dfrac{dx}{\sqrt{\pm(xx_1)(xx_2)(xx_3)(xx_4)}} = \dfrac{2\,dz}{\sqrt{(z^2-\alpha)(z^2-\beta)}}.}

Wenn nun die vier Wurzeln von $f(x)$ reell sind, so gibt es, wie wir im § 2 gesehen haben, acht Arten, diese Wurzeln den Zeichen $x_1$, $x_2$, $x_3$, $x_4$ so zuzuordnen, daß $a$ und $b$ entgegengesetzte Zeichen haben, und wenn also $\pm a$, $\mp b$ positiv sind, so werden $\sqrt{\alpha}$, $\sqrt{\beta}$ reell, also $\alpha$, $\beta$ positiv, und nach (7) wird dann auch $z$ reell.

Sind zweitens $x_1$, $x_2$ reell, $x_3$, $x_4$ konjugiert imaginär, so sind $a$ und $-b$ konjugiert imaginär, also wird, wenn die Vorzeichen der Quadratwurzeln $\sqrt{\pm a}$, $\sqrt{\mp b}$ passend bestimmt werden, $\sqrt{\alpha}$ reell, $\sqrt{\beta}$ rein imaginär, also $\alpha$ positiv, $\beta$ negativ und $z$ reell [weil $(x_3x_1)(x_4x_1)$ und $(x_3x_2)(x_4x_2)$ als Produkte konjugiert imaginärer Größen positiv sind].

Sind endlich $x_1$, $x_2$ und $x_3$, $x_4$ zwei konjugiert imaginäre Paare, so sind $a$ und $-b$ beide positiv. Nehmen wir also in (6) die unteren Zeichen, so werden $\sqrt{\alpha}$, $\sqrt{\beta}$ rein imaginär, also $\alpha$ und $\beta$ negativ, und aus (7) ergibt sich für $z$ ein reeller Ausdruck.
